\section{Conclusion}

We return to our opening question: how much does benchmark contamination inflate language model scores? Our checkpoint-gradient methodology provides a framework for answering this question. Preliminary validation suggests contamination-inflation correlation may exist at moderate strength (Spearman $r \approx 0.3$), pending full experimental validation with actual Pythia checkpoint evaluations.

This moderate correlation has important implications. Contamination does systematically inflate benchmark scores, confirming concerns about evaluation validity. Yet the effect is not overwhelming---capability remains the primary driver of benchmark performance. This suggests benchmark evaluation remains meaningful, but would benefit from contamination-aware correction.

Our methodological contributions---checkpoint-gradient analysis and capability detrending---enable contamination-performance studies without requiring contamination-free baseline models. Looking forward, we envision contamination-adjusted benchmark scores as standard practice, with our transfer function framework providing the foundation.

The path from contamination detection to contamination correction is now open. Our work takes the first step, demonstrating that contamination impact is not merely present but predictable.
